\begin{mistake}{2026-04-09}{28}

\begin{problemcontent}
求极限 \( \lim_{x \to \infty} \left( \sqrt[3]{(x+a)(x+b)(x+c)} - x \right) \)。
\end{problemcontent}

\begin{solutioncontent}
使用泰勒公式展开法。

首先将根式内部的多项式按降幂展开，并提取出最高次项 \( x^3 \)：
\( (x+a)(x+b)(x+c) = x^3 + (a+b+c)x^2 + (ab+bc+ca)x + abc \)。

原式可化为：
\[
\lim_{x \to \infty} x \left[ \left(1 + \frac{a+b+c}{x} + O\left(\frac{1}{x^2}\right)\right)^{1/3} - 1 \right]
\]

利用广义二项式展开公式 \( (1+u)^{\alpha} = 1 + \alpha u + o(u) \)（当 \( u \to 0 \) 时），此处 \( u = \frac{a+b+c}{x} + O(\frac{1}{x^2}) \)，\( \alpha = \frac{1}{3} \)：
\[
\left(1 + \frac{a+b+c}{x} + O\left(\frac{1}{x^2}\right)\right)^{1/3} - 1 \approx \frac{1}{3} \left( \frac{a+b+c}{x} \right) + O\left(\frac{1}{x^2}\right)
\]

代回极限式：
\[
\lim_{x \to \infty} x \left[ \frac{a+b+c}{3x} + O\left(\frac{1}{x^2}\right) \right] = \frac{a+b+c}{3}
\]
\end{solutioncontent}

\mistakeknowledge{
\begin{itemize}
 \item \textbf{核心公式}：广义二项式定理 \( (1+x)^\alpha = 1 + \alpha x + o(x) \) 的渐近展开应用。
 \item \textbf{易错点}：强行使用代数有理化公式 \( A^3-B^3 \) 进行展开，计算量极大极易出错；提取 \( x \) 的过程中遗漏指数 \( \frac{1}{3} \)，导致系数算错。
 \item \textbf{关联知识}：“抓大头”思想在无穷极限处的本质体现，利用提取主元法化归为 \( u \to 0 \) 的麦克劳林展开。
\end{itemize}}

\end{mistake}
